\section*{NeurIPS Paper Checklist}
% Draft answers; check the question list against the NeurIPS 2027 call for papers when it is published (G9_PLAN O5).
\begin{enumerate}
\item \textbf{Claims.} Do the main claims in the abstract and introduction accurately reflect the paper's contributions and
scope? \textbf{[Yes]} Every claim is traced to a theorem or a pre-registered test (\cref{sec:theory,sec:exp}); exploratory
analyses are labelled as such; the scope (LLP with bags of 8 and PLL; dependence on the TFM) and the two failed criteria
of Study~3 are stated in the abstract, \cref{sec:s3} and \cref{sec:limitations}.
\item \textbf{Limitations.} \textbf{[Yes]} \Cref{sec:limitations}, including the unmet hypotheses of Study~1, the failed criteria of
Study~3, the scope of each guarantee, and the exploratory status of the comparison with the tuned learner in Study~3.
\item \textbf{Theory assumptions and proofs.} \textbf{[Yes]} Assumptions (A0)--(A3) in \cref{sec:setup} and \cref{ass:H}; complete proofs in
\cref{app:proofs}, numerically checked in \cref{app:rq1,app:blockbias,app:dynamics}.
\item \textbf{Experimental result reproducibility.} \textbf{[Yes]} Protocols, splits, weak-label generators, tuning spaces
and all settings are specified (\cref{sec:design}, \cref{app:data}); every number is generated by scripts from raw
per-cell results.
\item \textbf{Open access to data and code.} \textbf{[Yes]} Public datasets; anonymised code and raw results
(\cref{app:statements}).
\item \textbf{Experimental setting/details.} \textbf{[Yes]} \Cref{sec:method,sec:design}, \cref{app:data}.
\item \textbf{Experiment statistical significance.} \textbf{[Yes]} Wilcoxon signed-rank tests with Holm correction,
bootstrap confidence intervals, win/tie/loss counts, Bayesian signed-rank tests, and a registered sensitivity analysis.
\item \textbf{Experiments compute resources.} \textbf{[Yes]} \Cref{sec:ablation} (cost per run) and \cref{app:statements}.
\item \textbf{Code of ethics.} \textbf{[Yes]} The research uses public data and conforms to the NeurIPS Code of Ethics.
\item \textbf{Broader impacts.} \textbf{[Yes]} Learning from weak labels lowers the cost of annotation. Learning from label
proportions also recovers individual-level predictions from data released only in aggregate, which can weaken the
privacy that aggregation was meant to provide (\cref{app:statements}).
\item \textbf{Safeguards.} \textbf{[NA]} No model or data with high risk of misuse is released.
\item \textbf{Licenses for existing assets.} \pending{licences of the TabICL~v2, TabPFN-3 and TabDPT checkpoints, each OpenML dataset and
the PALM partial-label data not yet confirmed} Draft answer, to be confirmed asset by asset: TabPFN~v2 weights: Prior Labs
License (Apache~2.0 with an attribution requirement); TabPFN-3 weights: the publisher's non-commercial licence; TabDPT
code and weights: Apache~2.0 as listed by the publisher; TabICL code: BSD-3-Clause; datasets: licences as listed on OpenML;
baseline code without a licence is loaded at run time and not redistributed.
\item \textbf{New assets.} \textbf{[Yes]} Code and raw results, documented in the repository.
\item \textbf{Crowdsourcing and research with human subjects.} \textbf{[NA]}
\item \textbf{IRB approvals.} \textbf{[NA]}
\item \textbf{Declaration of LLM usage.} \textbf{[Yes]} \Cref{app:statements}.
\end{enumerate}
